\section{Introduction}
\label{sec:introduction}

The Border Gateway Protocol (BGP) is the foundational protocol that enables inter‐domain routing across tens of thousands of Autonomous Systems (ASes)~\cite{rfc4271}. 


The inter‐domain BGP validation and filtering, introduced in 1995, have relied on the Internet Routing Registry (IRR)~\cite{DuIRATSc23}, a public database that allows operators to publish IP‐prefix/ASN pairs, with each prefix authorized to originate from a given ASN. Nevertheless, IRRs provide neither strong publisher authentication nor cryptographic guarantees of prefix ownership, leading to persistent risks of misinformation or misconfiguration~\cite{MahajanWA02,KangLRKC24IRRedicator,DuIRATSc23}. Moreover, conflicting data across multiple IRRs creates loopholes that adversaries can exploit to bypass route‐monitoring filters~\cite{BirgeLeeAR25}.
The absence of a native authentication mechanism exposes BGP to a wide range of attacks, including hijacking~\cite {prefixHijacking} and leakage~\cite {McDanielSS21BGPLeak}, resulting in severe disruptions, such as traffic blackholing, interception, or global outages~\cite{butler2009survey}. 

%Recent work demonstrates that even RPKI-protected networks remain vulnerable to update-message flooding attacks that induce route oscillations that can destabilize large portions of the Internet~\cite{stoger2025vortex}.


While RPKI \cite{RPKI} was standardized in 2012 to provide a cryptographic foundation for origin validation, its effectiveness is strictly limited by its binary nature (\texttt{VALID} and \texttt{INVALID}) and incomplete global adoption. Especially, the high barrier to entry and lack of immediate economic incentives have led to a stagnation in deployment, with only 27\% of ASes currently enforcing validation. Consequently, the Internet remains a fragmented landscape where the cryptographic root of trust is frequently broken at the boundary between RPKI-enforcing and non-participating networks, necessitating a more inclusive, reputation-based extension to the current security architecture.



Although existing research has criticized RPKI's centralized trust model and its complex fetching and caching mechanisms~\cite{HlavacekSVW23BGPAdopt}, and recent systematization reveals that 56\% of RPKI validators contain implementation vulnerabilities exploitable for Denial of Service (DoS) threats~\cite{stoger2025vortex, mirdita2025sok}, \rj{raising \textbf{single-point failure} concerns. To address the systemic vulnerabilities that undermine the effectiveness of RPKI in practice, this work investigates three core research questions aimed at reinforcing ROA integrity in real-world operations:}

\noindent\textbf{(RQ1) Bridging the Deployment Gap.} \rj{How can the cryptographic root of trust from the RPKI ecosystem be transitively extended to provide verifiable security for the 63\% of the Internet that remains outside the RPKI periphery?}



Despite significant progress in prefix signing, where 47.7\% of IPv4 and 50.45\% of IPv6 prefixes are now covered by valid ROAs \cite{HlavacekSVW23BGPAdopt}, a critical enforcement gap remains. The majority of the Internet still operates without active RPKI-based validation on BGP sessions, meaning that even when cryptographic records exist, they are frequently ignored during route selection \yiwen{please add citations here to support your claim}. This leaves the remainder of the ecosystem unable to programmatically verify incoming announcements, rendering the existing cryptographic root ineffective at the network edge. Furthermore, \textbf{hybrid approaches} that bridge this gap by combining \textbf{RPKI} with legacy \textbf{IRR} filtering \cite{KangLRKC24IRRedicator} continue to suffer from the weak authentication inherent in IRR databases, thereby \textit{introducing additional attack vectors rather than resolving the underlying trust deficit}.

\noindent \textbf{(RQ2) Establishing an Audit Trail.} \rj{How can a decentralized and immutable ledger be designed to transform transient BGP routing observations into a longitudinal and scalable audit trail of AS behavior?}

Existing \textbf{blockchain-based BGP security solutions}, such as RouteChain~\cite{saad2022routechain} and ISRchain~\cite{chen2020isrchain}, fail to provide a reliable audit trail due to two structural deficiencies. \rj{First, they \textit{lack topological alignment}.} RouteChain groups ASes by geographical proximity rather than functional peering relationships, resulting in an audit record that misinterprets the path of trust. \rj{Second, they suffer from a \textit{data integrity gap}.} By relying on passive logging without cryptographically anchoring observers to the RPKI root, these systems cannot verify the authority of the entity reporting a routing event. Consequently, their ledgers remain collections of unverified, fragmented snapshots rather than a continuous, authoritative history capable of supporting longitudinal behavioral analysis.

\noindent \textbf{(RQ3) Enabling Security Enforcement.} \rj{What is the missing link that can effectively translate decentralized behavioral audits into active and secure routing enforcement?}

The RPKI framework (RFC 6480~\cite{RPKI}) provides transient origin validation but lacks a standardized mechanism for longitudinal state retention. This creates an \rj{\textbf{enforcement blind spot}} without a verifiable record of past behavior; thereby, operators cannot differentiate between isolated misconfigurations and persistent malicious hijacking. Consequently, there is a critical need for an RPKI-based architecture that: 1) transforms transient observations into a persistent, verifiable audit trail, and 2) translates this behavioral history into \rj{automated and trust-based enforcement policies} that extend security guarantees even to non-RPKI networks.

\rj{\textbf{The design goal of BGP-Sentry is to \textit{illuminate the global AS topology} by extending the reach of cryptographic trust to the previously opaque regions of the Internet. Currently, the inter-domain landscape is a fragmented ``isolated world,'' split between the RPKI-protected core and a vast unverifiable periphery. To tackle the challenge, we propose to bridge this paradox, expanding the ``trust'' of the Central RPKI to unify these disconnected domains into a single, accountable routing ecosystem.}}


\rj{The proposed BGP-Sentry is a \textbf{blockchain-based trust-expansion framework} that transforms \textit{RPKI-enabled ASes} into \textit{distributed observers of their neighbors}. By recording observed BGP behaviors on a shared ledger \textit{rooted in the foundational trust of the central RPKI}, these events are used to compute a longitudinal \textit{Operational Attestation Level (OAL)} for each AS (both RPKI-enabled and non-RPKI) through protocol-embedded algorithms executed during blockchain consensus.}

% The design goal of BGP-Sentry is to provide a robust system foundation capable of supporting diverse behavioral metrics to illuminate the inter-domain landscape. 
% To demonstrate this flexibility, the architecture is designed to ingest signals such as \textit{1) announcement consistency} to detect route flapping, \textit{2) reverse-path verification} to ensure reachability, and \textit{3) registry cross-checking} to align BGP data with external databases. 
% These metrics are synthesized into a longitudinal \textit{Operational Attestation Level (OAL)}, which moves beyond transient "scores" to reflect an AS's long-term protocol integrity. 

% As a conceptual showcase of this framework, ASes can be qualitatively classified into five tiers, e.g., \texttt{Highly} \texttt{Trusted}, \texttt{Trusted}, \texttt{Neutral}, \texttt{Suspicious}, and \texttt{Malicious}. While the precise algorithmic weightings and numerical thresholds for these tiers are designated as future work, this classification provides an immediate dual benefit: it establishes a permanent, verifiable audit trail of historical behavior and enables RPKI-compliant operators to make informed, risk-aware routing decisions when interfacing with non-RPKI domains.

%Lastly, to ensure the viability of such a distributed system, BGP Coins are introduced as a participation incentive. Rather than relying on altruism, these tokens motivate RPKI-enabled ASes to join the consortium as active observers, providing the collective monitoring power necessary to unify isolated routing worlds into a single, accountable ecosystem.


%Furthermore, BGP‑Sentry employs behavioral metrics, including announcement consistency, reverse‑path verification, registry cross‑checking, and endorsement alignment, to dynamically update trust scores. Finally, ASes can be conceptually classified into five tiers, e.g., Highly Trusted ($\geq$90), Trusted (70--89), Neutral (50--69), Suspicious (30--49), and Malicious ($<$30), providing both a historical audit trail among RPKI ASes and informed routing decisions for non-RPKI ASes. Lastly, to encourage participation, BGP Coins incentivize RPKI-enabled observers to maintain honest and accurate monitoring of their non-RPKI neighbors. Our contributions are as follows:

\noindent \textbf{(i) Resolving Isolation via Hierarchical Chain of Trust.} 
\rj{We resolve the cryptographic opacity and ``isolated world'' problem of non-RPKI networks by designing a multi-layered chain of trust that extends the cryptographic authority of the Central RPKI to the previously opaque non-RPKI regions. By anchoring blockchain identities to RPKI certificates, the framework establishes a verifiable provenance for all routing attestations, allowing RPKI-enabled ASes to act as ``trust proxies'' and enabling authoritative attribution and forensic accountability for 63\% of the Internet currently lacking native cryptographic protection.}

\noindent \textbf{(ii) Scaling Up Identity-based Blockchain Consensus.} 
\rj{We resolve the scalability bottlenecks of identity-based blockchains, specifically $O(n^2)$ communication complexity and recursive validation overhead, by introducing a \textit{Data Plane-aware Proof of Population} (DP-PoP) consensus anchored to the RPKI hierarchy. By restricting the validator set to verified and neighbour ASes within a few hops, BGP-Sentry achieves sub-linear overhead growth and consensus latencies under 3.2s. This foundation supports an incentivized monitoring ecosystem via \textit{BGP Coins} and \textit{(OAL)}, effectively extending the RPKI root of trust to 63\% of non-RPKI networks to create a proactive, self-defending routing infrastructure.}

\noindent \textbf{(iii) Longitudinal Behavioral Attestation and Tiered Trust.} 
\rj{We propose a robust framework for inter-domain accountability by synthesizing diverse BGP signals, including announcement consistency, reverse-path verification, and registry cross-checking, into a longitudinal \textit{Operational Attestation Level (OAL)}. Unlike transient reputation scores, the OAL captures long-term protocol integrity through protocol-embedded algorithms executed during blockchain consensus. We demonstrate the utility of this framework through a qualitative five-tier trust classification (from \texttt{Highly Trusted} to \texttt{Malicious}), providing RPKI-compliant operators with an immutable audit trail and a deterministic basis for risk-aware routing decisions when interfacing with the 63\% of non-RPKI domains.}

\noindent \textbf{(iv) Incentive-Compatible Monitoring via BGP Coins.} 
\rj{We ensure the long-term viability of the decentralized framework by introducing \textit{BGP Coins}, a participation incentive that replaces operator altruism with tangible rewards for active observation. By aligning the economic interests of RPKI-enabled ASes with global security goals, this mechanism fosters the collective monitoring power necessary to bridge isolated routing domains into a unified, accountable, and self-defending ecosystem.}

%\noindent\textbf{(i) Expanding Trust-based Routing to Non-RPKI ASes.} By turning RPKI‑enabled ASes into voluntary, distributed BGP observers, we mitigate the real-world challenge of deploying RPKI across the entire network. The proposed blockchain with protocol-embedded trust score computing expands the RPKI root trust to non-RPKI networks and allows informed routing decision-making for RPKI-enabled ASes interfacing unknown ASes, enhancing overall network robustness.

%\noindent\textbf{(ii) Enabling Behavioral Monitoring and Observer Incentives.} The RPKI-based blockchain onboarding process preserves the root of trust; thus, BGP-Sentry minimizes the risk of misbehavior of blockchain participants (RPKI-enabled ASes) and achieves reliable trust score assessment for non‑RPKI ASes. The BGP Coin incentive system motivates RPKI observers to provide honest and accurate monitoring, while trust scores create reputation-based consequences for non-RPKI ASes, encouraging improved routing behavior.

%\noindent\textbf{(iii) Establishing Auditable Trust Rating through Scalable Consensus} BGP-Sentry employs a Proof of Population (PoP) consensus mechanism (i.e., one vote per node), enabling all RPKI observers to agree on behavioral trust ratings for non-RPKI ASes. These consensus-driven ratings are recorded on the blockchain, providing immutable and auditable trust assessments. Evaluation demonstrates scalable operation across large networks with negligible overhead and sub-linear growth.

%The rest of this paper is organized as follows: we first outline the motivation for our system design in section~\ref{sec:motivation}. Then, in section~\ref{sec:architecture}, we describe system architecture and workflow. Next, in section~\ref{sec:evaluation}, we evaluate the proposed system in terms of overhead and security enforcement, with extended discussions located in section~\ref{sec:discussion}. Lastly, we discuss related works and conclude our work in sections~\ref{sec:relatedWork} and~\ref{sec:conclusion}, respectively.